% META: {"id":"1NSI-TC-CO-005","chapitre":"1NSI-TYPES-CONSTRUITS","type_objet":"corrige","capacites":["1NSI-TYPES-CONSTRUITS-C1"],"mode_creation":"ex_nihilo","sources_inspiration":[],"status":"approved","exercice_ref":"1NSI-TC-EX-005","origin":"needs_review","status_history":["needs_review","audited","accepted","approved"],"release_acceptance":"RELEASE_OWNER_BATCH_ACCEPTANCE","acceptance_closure_digest":"sha256:9b3ccf9a81c5520fbb7e03b7b2d3e7bf2057a02834b6d908bf3be5f49f3b553f","acceptance_receipt_digest":"sha256:4fad86f0282e0fa4e0419cca1ad6379ae7af5a280c04a4a90880f90a1c044242"}
\begin{corrige}{1NSI-TC-EX-005}
\begin{enumerate}
  \item Le code affiche \lstinline{Paris} puis \lstinline{Lyon}. En effet, seules ces deux villes ont une température strictement supérieure à 30 (32.5 et 35.1). Lille (28.7) ne satisfait pas la condition.

  \item Chaque élément de \lstinline{releves} est un \textbf{tuple} de la forme \lstinline{(str, float)}.

  \item \lstinline{releves[1][0]} vaut \lstinline{"Lyon"} : c'est le premier élément (indice 0) du deuxième tuple (indice 1) de la liste.
\end{enumerate}

Vérification : les résultats ont été confirmés par exécution.
\end{corrige}
